% =============================================================
\section{Related Work}\label{sec:related}

Our contribution sits at the confluence of three literatures that have, to
date, evolved in isolation: (i)~differentiable parametric surface modelling
for CAD, (ii)~surface-bound or surface-extracted Gaussian splatting, and
(iii)~intrinsic texture representations on manifolds. We review each in turn
and identify the empty intersection that motivates this work
(\autoref{tab:landscape}).

% -------------------------------------------------------------
\subsection{Differentiable Parametric Surface Modelling}\label{sec:rw:diffcad}

\paragraph{Differentiable NURBS and B-splines.}
NURBS-Diff~\cite{nurbsdiff} introduced a differentiable programming module for
NURBS curves and surfaces integrated into PyTorch, demonstrating unsupervised
point-cloud reconstruction and analysis-constrained surface fitting. The
Diff-NURBS library~\cite{diffnurbs} provides a maintained, pure-PyTorch
NURBS implementation with closed-form surface derivatives and normals under
an Apache-2.0 licence, serving as the geometric backbone of our method.
VALVEFIT~\cite{valvefit} applies GPU-accelerated differentiable B-spline
surface fitting to patient-specific heart-valve modelling, while
NURBS-PyTorch~\cite{nurbspytorch} and TorchNurbs~\cite{torchnurbs} are earlier
differentiable NURBS implementations now largely unmaintained. A common
limitation of this entire family is the absence of a rendering pipeline: the
fitting signal is a point cloud or a prescribed loss, never multi-view
images, so the geometric and appearance channels have not been coupled
end-to-end with a differentiable renderer.

\paragraph{Differentiable subdivision surfaces.}
As an alternative to parametric splines, differentiable subdivision surface
fitting~\cite{diffsubsurf} fits Loop subdivision surfaces to point clouds via
an implicit moving-least-squares target coupled with a differentiable
rasteriser. Subdivision handles arbitrary topology more gracefully than
tensor-product B-splines, but yields neither a compact control mesh with an
analytic metric nor a CAD-native output. Our choice of B-splines sacrifices
arbitrary topology (we restrict the present work to single-patch CAD-class
objects) in exchange for a closed-form first fundamental form and direct
compatibility with manufacturing pipelines.

\paragraph{CAD reconstruction and reverse engineering.}
A complementary line reconstructs boundary representations (B-reps) from
unstructured input. DualBrep~\cite{dualbrep} (SIGGRAPH~2026) unifies B-rep
geometry and topology in dual SDF/UDF scalar fields and extracts explicit
B-rep models via flow matching from point clouds. A2Z-10M+~\cite{a2z} and
PS-CAD~\cite{pscad} learn CAD modelling sequences from large annotated
datasets. These methods operate on point clouds or modelling commands and do
not employ differentiable rendering of the parametric surface itself; they
are orthogonal to our image-based, rendering-driven formulation and could
supply initial geometry for our joint optimisation.

\paragraph{NURBS splatting in the plane.}
Closest in spirit to our rendering formulation is NURBS
Splatting~\cite{nurbssplatting} (ECCV~2026), which renders planar rational
spline curves as continuous Gaussian fields by sampling Gaussians along the
curve parameter domain, reformulating vector-graphics rendering as a smooth
Gaussian-accumulation process with stable gradients. The framework elegantly
handles rational weights, non-uniform knots, and closed-region filling.
Crucially, however, it is confined to \emph{two-dimensional} planar curves:
there is no 3D surface, no multi-view consistency, and no joint
geometry--appearance optimisation. Our work can be read as the
three-dimensional, image-driven generalisation of this idea to B-spline
\emph{surfaces} carrying intrinsic heat-kernel appearance.

% -------------------------------------------------------------
\subsection{Surface-Bound Gaussian Splatting}\label{sec:rw:bindinggs}

A substantial body of work tames the geometric ambiguity of 3DGS by binding
Gaussian primitives to an explicit surface. SuGaR~\cite{sugar} regularises
Gaussians to align with the underlying surface and optionally binds them to a
Poisson-reconstructed mesh for joint optimisation and editing.
OMeGa~\cite{omega} jointly optimises an explicit mesh and 2D splats whose
spatial attributes are expressed in the mesh frame, improving indoor-scene
geometry. MeshGS~\cite{meshgs} distinguishes tightly- and loosely-bound
splats to align some Gaussians to geometry while letting others absorb mesh
artefacts. GauSTAR~\cite{gaustar} binds Gaussians to mesh faces for dynamic
surface tracking and reconstruction, unbinding them where topology changes.
Manifold-GS~\cite{manifoldgs} represents surface-like Gaussians as discrete
unoriented varifolds and exports certified open surface patches.
EGG-Fusion~\cite{eggfusion} performs real-time geometry-aware Gaussian-surfel
mapping for SLAM. Earlier foundations include 2D Gaussian
Splatting~\cite{2dgaussians} and geometry-grounded
primitives~\cite{geosplatting}.

A unifying property of \emph{all} the above is the surface carrier: a
\emph{discrete triangular mesh}. The mesh is $C^0$ across edges, possesses no
analytic metric tensor, and is not a CAD-native representation. Consequently,
kernel covariance must be inferred from discrete geometry, and the output
cannot be fed directly to a manufacturing pipeline without reconstruction.
Our formulation replaces the mesh with a differentiable B-spline surface,
yielding a closed-form metric, $C^2$ continuity, and native CAD
compatibility---at the cost of restricting the present scope to
single-patch, CAD-class topology.

% -------------------------------------------------------------
\subsection{Intrinsic Texture and Appearance on Manifolds}\label{sec:rw:hktex}

Heat Kernel Textures~\cite{hktex} (ECCV~2026) represent appearance on a
triangular mesh as a set of anisotropic heat kernels---the geodesic
counterparts of Gaussians---defined directly on the surface. Kernel
positions, optimisation, pruning, and densification all operate in
barycentric coordinates, avoiding UV seams, distortion, wasted atlas space,
and uneven texel resolution. The representation is fitted from an existing
texture or from multi-view images and integrates with physically based
rendering.

The computation of HKTex, however, is intrinsically mesh-discrete: it relies
on a discrete Laplacian eigen-decomposition and a $k$-nearest-neighbour heat
propagation to approximate the geodesic quantities that define each kernel.
Our contribution is a strict generalisation: on a B-spline surface the first
fundamental form is available in closed form and is differentiable in the
control points, so each anisotropic heat kernel reduces to a Gaussian on the
$(u,v)$ domain with covariance $g^{-1}$, and the entire eigen-decomposition
and KNN machinery is eliminated. This simultaneously simplifies the
formulation, makes it differentiable with respect to the geometry, and yields
a representation that lives on a smooth rather than a discrete manifold.

Related appearance-on-surface work includes TextureSplat~\cite{texturesplat}
and SVG-IR~\cite{svgir}, which attach per-primitive spatially-varying
materials to Gaussians, and the broader inverse-rendering line of
GS-IR~\cite{gsir}, Relightable~3D~Gaussians~\cite{relightable3dgaussians},
DeferredGS~\cite{deferredgs}, and their successors~\cite{gigs,gusir,rtrgs}.
These methods model appearance as per-primitive attributes or small MLPs
rather than intrinsic manifold kernels, and they remain mesh- or
primitive-bound.

% -------------------------------------------------------------
\subsection{The Empty Intersection}\label{sec:rw:gap}

\begin{table}[t]
  \caption{Landscape of adjacent literatures. The cell marked
  \textbf{(ours)} is, to our knowledge, unoccupied.}
  \label{tab:landscape}
  \small
  \begin{tabular}{lccc}
    \toprule
     & Gaussian rendering & Diff.~rasterisation & Generative / flow \\
    \midrule
    2D planar curves     & NURBS Splatting~\cite{nurbssplatting} & --- & --- \\
    3D triangular mesh   & SuGaR / OMeGa / GauSTAR~\cite{sugar,omega,gaustar} & DiffSubSurf~\cite{diffsubsurf} & --- \\
    3D parametric surface & \textbf{(ours)} & --- & DualBrep~\cite{dualbrep} \\
    3D B-rep             & --- & --- & DualBrep~\cite{dualbrep} \\
    \bottomrule
  \end{tabular}
\end{table}

\autoref{tab:landscape} summarises the landscape. NURBS
Splatting~\cite{nurbssplatting} establishes the spline-as-Gaussian-field
rendering paradigm but only for 2D curves. Surface-bound
Gaussian splatting~\cite{sugar,omega,gaustar,meshgs,manifoldgs} binds
Gaussians to surfaces but only to discrete meshes. HKTex~\cite{hktex}
defines intrinsic Gaussian appearance but only on a discrete mesh with a
fixed carrier. Differentiable CAD~\cite{nurbsdiff,diffnurbs,valvefit} and
generative B-rep reconstruction~\cite{dualbrep} model parametric geometry but
do not render it with Gaussians and do not optimise appearance jointly from
images. The intersection---a differentiable B-spline \emph{surface} rendered
as intrinsic Gaussian appearance and optimised jointly from multi-view
images---is empty, and this paper occupies it.
